\newpage
\section*{NeurIPS Paper Checklist}

\begin{enumerate}

\item {\bf Claims}
    \item[] Question: Do the main claims made in the abstract and introduction accurately reflect the paper's contributions and scope?
    \item[] Answer: \answerYes{}
    \item[] Justification: The abstract and introduction state contributions C1--C6 and report the headline numbers ($53/60$ on the historical bug corpus; $32/34$ on the modern Pytea-fragment subset; $0$ unconditional Refuted-Proof on the $488$-block applicability survey; $11/11$ Lean lemmas closed sorry-free); the scope of the soundness theorem and the unmechanised parts of the implementation are stated explicitly in \Cref{sec:eval-lean,sec:limconc}.

\item {\bf Limitations}
    \item[] Question: Does the paper discuss the limitations of the work performed by the authors?
    \item[] Answer: \answerYes{}
    \item[] Justification: \Cref{sec:limconc} discusses the supported language fragment, the constructor-bound integer-attribute envelope class of remaining silent verifieds, the gap between the Lean-audited operator-rule table and the un-mechanised Python analyser, the first-order grad-flag lattice's $\le 12\%$ silent-misclassification regime on parameter-sharing under renamed attributes, and the small-$N$ statistical limits of the post-freeze unfiltered evaluation.

\item {\bf Theory Assumptions and Proofs}
    \item[] Question: For each theoretical result, does the paper provide the full set of assumptions and a complete (and correct) proof?
    \item[] Answer: \answerYes{}
    \item[] Justification: Theorems~\ref{thm:soundness} (soundness), \ref{thm:preservation}--\ref{thm:progress} (preservation/progress for the calculus), the assume/guarantee composition theorem, and the Dynamo-correspondence lemma (\Cref{thm:dynamo-corr}) are stated with all assumptions and proved in the appendix; \textsc{Library-Warn} and \textsc{Abstain} are explicitly excluded from the soundness claim. The Lean~$4$ operator-rule audit is sorry-free at the rule definitions, with the soundness lemma scope (28 of 79 handlers) tabulated in \Cref{tab:handler-soundness}.

\item {\bf Experimental Result Reproducibility}
    \item[] Question: Does the paper fully disclose all the information needed to reproduce the main experimental results of the paper to the extent that it affects the main claims and/or conclusions of the paper?
    \item[] Answer: \answerYes{}
    \item[] Justification: All evaluation numbers are produced by deterministic single-CPU driver scripts whose inputs (corpora, freeze hashes, pre-registered queries) are committed; the released artefact contains the verifier source, the Lean development, the benchmark drivers, and the JSON inputs that produce every reported table.

\item {\bf Open access to data and code}
    \item[] Question: Does the paper provide open access to the data and code, with sufficient instructions to faithfully reproduce the main experimental results, as described in supplemental material?
    \item[] Answer: \answerYes{}
    \item[] Justification: An anonymised supplementary archive accompanies the submission; it contains the verifier source, the Lean development, all benchmark drivers, the committed JSON corpora, and a top-level README giving the exact commands needed to reproduce every table and figure.

\item {\bf Experimental Setting/Details}
    \item[] Question: Does the paper specify all the training and test details necessary to understand the results?
    \item[] Answer: \answerYes{}
    \item[] Justification: There is no model training. \Cref{sec:eval} and the appendix specify the analysed modules, the input shape envelopes, the analyser configuration, the catalogue freeze commit, the modern-subset construction predicate, the post-freeze pre-registered query date, and the matched-pair test used for the Pytea comparison.

\item {\bf Experiment Statistical Significance}
    \item[] Question: Does the paper report error bars suitably and correctly defined or other appropriate information about the statistical significance of the experiments?
    \item[] Answer: \answerYes{}
    \item[] Justification: Wilson $95\%$ confidence intervals are reported for the bug-corpus and post-freeze catch rates; a McNemar exact two-sided test and paired-bootstrap interval are reported for the modern-subset Pytea head-to-head; the post-freeze unfiltered $N{=}15$ comparison reports Fisher exact $p$-values and is explicitly described as not separable at $\alpha{=}0.05$.

\item {\bf Experiments Compute Resources}
    \item[] Question: For each experiment, does the paper provide sufficient information on the computer resources needed to reproduce the experiments?
    \item[] Answer: \answerYes{}
    \item[] Justification: The full benchmark suite runs single-threaded on one CPU core in under a few minutes; the Lean development builds in under $10$ seconds via \texttt{lake build}; no GPU is used and no preliminary or failed runs required additional compute.

\item {\bf Code Of Ethics}
    \item[] Question: Does the research conducted in the paper conform, in every respect, with the NeurIPS Code of Ethics?
    \item[] Answer: \answerYes{}
    \item[] Justification: The work is a static program-analysis tool involving no human subjects, no data collection, no model training, and no released model; it conforms to the NeurIPS Code of Ethics.

\item {\bf Broader Impacts}
    \item[] Question: Does the paper discuss both potential positive societal impacts and negative societal impacts of the work performed?
    \item[] Answer: \answerYes{}
    \item[] Justification: \Cref{sec:limconc} notes that catching shape and gradient bugs statically reduces wasted GPU time and developer effort; no negative societal impact is foreseeable since the released artefact is a developer-facing static checker rather than a generative or decision-making system.

\item {\bf Safeguards}
    \item[] Question: Does the paper describe safeguards that have been put in place for responsible release of data or models that have a high risk for misuse?
    \item[] Answer: \answerNA{}
    \item[] Justification: The paper releases only a static verifier and benchmark drivers; no pretrained model, generative model, or scraped dataset is released, so no misuse-prevention safeguards are required.

\item {\bf Licenses for existing assets}
    \item[] Question: Are the creators or original owners of assets used in the paper properly credited and are the license and terms of use explicitly mentioned and properly respected?
    \item[] Answer: \answerYes{}
    \item[] Justification: The PyTorch modules used as analysis targets (\textsc{torchvision}, BSD-3; \textsc{timm}, Apache-2.0; \textsc{transformers}, Apache-2.0) are cited in the paper, used solely as analysis inputs, not redistributed, and their licenses are respected; PyTorch itself ($2.9.1$, BSD-style) and Pytea (MIT) are similarly cited.

\item {\bf New Assets}
    \item[] Question: Are new assets introduced in the paper well documented and is the documentation provided alongside the assets?
    \item[] Answer: \answerYes{}
    \item[] Justification: The released artefacts are the verifier source, the Lean operator-rule development, and the benchmark drivers, all under Apache-2.0; they are documented by an anonymised top-level README and inline docstrings shipped with the supplementary material. No model checkpoints or datasets are released.

\item {\bf Crowdsourcing and Research with Human Subjects}
    \item[] Question: For crowdsourcing experiments and research with human subjects, does the paper include the full text of instructions given to participants and screenshots, if applicable, as well as details about compensation (if any)?
    \item[] Answer: \answerNA{}
    \item[] Justification: The paper does not involve crowdsourcing or any research with human subjects.

\item {\bf Institutional Review Board (IRB) Approvals or Equivalent for Research with Human Subjects}
    \item[] Question: Does the paper describe potential risks incurred by study participants, whether such risks were disclosed to the subjects, and whether IRB approvals were obtained?
    \item[] Answer: \answerNA{}
    \item[] Justification: The paper does not involve human subjects, so no IRB approval or equivalent is required.

\end{enumerate}
